% =====================================================================
%  Aula 8 — Seleção Adversa e Risco Moral (MPE 2026/32, Insper)
%  TRÊS EXERCÍCIOS RESOLVIDOS PARA LOUSA (aula presencial de revisão)
%  Calibre: piso Nicholson & Snyder 12e Cap. 18; teto Jehle-Reny 3e.
%
%  Compilar com: pdflatex exercicios-lousa-aula8.tex   (ou lualatex)
%
% ---------------------------------------------------------------------
%  DECISÃO: TRÊS exercícios (fechada pelo professor em 2026-06-10).
%  --------------------------------------------------------------
%  Os três pilares da Aula 8: Akerlof (hidden type) / Holmström (hidden
%  action) / Rothschild-Stiglitz (separador). O 4º (pooling regulado
%  ANS/ACA) foi descartado — reusava o avaliativo 6 da plataforma e
%  estourava o buffer de tempo. Não há mais esqueleto neste arquivo.
%
% ---------------------------------------------------------------------
%  DOIS PDFs A PARTIR DESTA FONTE ÚNICA (padrão \ifsolucao do projeto):
%  -------------------------------------------------------------------
%   - Versão LOUSA (com gabarito): compile ESTE arquivo direto.
%       pdflatex exercicios-lousa-aula8.tex
%   - Versão CADERNO DO ALUNO (só enunciados): compile o wrapper.
%       pdflatex caderno-aluno-aula8.tex
%  Sem duplicar conteúdo: o wrapper define \modoaluno e dá \input aqui.
%
% ---------------------------------------------------------------------
%  ORÇAMENTO DE TEMPO DE LOUSA (professor escrevendo + explicando)
%  ----------------------------------------------------------------
%   Ex 1 — Akerlof (unraveling parcial vs. colapso) .......  ~40 min
%   Ex 2 — Holmström (first-best vs. second-best) .........  ~55 min
%   Ex 3 — Rothschild-Stiglitz (separador) ................  ~45 min
%   ------------------------------------------------------------------
%   CONTEÚDO PURO ......................................... ~140 min
%   BUFFER (perguntas, idas ao quadro, respiro) ...........  ~20 min
%   ------------------------------------------------------------------
%   TOTAL ALVO ............................................  160 min  (2h40)
%
%  Calibragem honesta: 140 de conteúdo + 20 de buffer é APERTADO mas
%  cabe. Variável de ajuste em tempo real: se Ex 1 derrapar, cortar a
%  parte (c) do Ex 2 (equivalência sob neutralidade) — é a única peça
%  conceitualmente destacável sem ferir o arco. Se sobrar tempo, abrir
%  a discussão ANS/ACA do comentário do Ex 3 ou puxar o 4º exercício.
%
% ---------------------------------------------------------------------
%  MAPA CONCEITO -> EXERCÍCIO (o que cada um RESUME da Aula 8)
%  ----------------------------------------------------------
%   Ex 1  : hidden type; Akerlof (1970) unraveling; média condicional;
%           colapso (multiplicativo) vs. mercado parcial (aditivo);
%           perda de bem-estar por seleção adversa.
%   Ex 2  : hidden action; principal-agente Holmström (1979);
%           first-best (salário fixo, P absorve risco) vs. second-best
%           (IR+IC binding, w_H>w_L); prêmio de risco moral;
%           equivalência sob neutralidade do agente (contrato franchise).
%   Ex 3  : hidden type em seguro; Rothschild-Stiglitz (1976) separador;
%           single-crossing; H cobertura completa / L distorcido;
%           IC do H binding; (menção) inexistência sob lambda baixo.
%
%  O QUE FICOU DE FORA DE PROPÓSITO (não cabe em 160 min com qualidade):
%   - Derivação FORMAL da inexistência R-S (precisa de teoria de jogos
%     contínua) -> só menção + intuição "pooling não é Nash" no Ex 3.
%   - Mirrlees (1971) e Grossman-Hart (1983) -> são panorâmicos no
%     Bloco 4 da aula; não rendem exercício de lousa enxuto.
%   - Sinalização (Spence) e matching -> são da AULA 9; não antecipar.
% =====================================================================

\documentclass[11pt,a4paper]{article}

\usepackage[utf8]{inputenc}
\usepackage[T1]{fontenc}
\usepackage[brazilian]{babel}
\usepackage{amsmath,amssymb}
\usepackage{mathtools}
\usepackage[margin=2.4cm]{geometry}
\usepackage{enumitem}
\usepackage{xcolor}
\usepackage[most]{tcolorbox}
\usepackage{booktabs}

% --- Cor institucional Insper ---
\definecolor{insper}{HTML}{C8102E}
\definecolor{mundo}{HTML}{1B3A5C}
\definecolor{amarelo}{HTML}{C79A00}  % calibre intermediário (N&S)  [🟡]
\definecolor{verm}{HTML}{C8102E}     % calibre avançado (Jehle-Reny) [🔴]
\definecolor{verde}{HTML}{2E7D32}    % calibre consolidação          [🟢]

% --- Marcadores de calibre (substituem os emojis 🟢🟡🔴 sob pdflatex) ---
\newcommand{\calY}{\textbf{\textcolor{amarelo}{[N]}}} % intermediário N&S
\newcommand{\calR}{\textbf{\textcolor{verm}{[J-R]}}}  % avançado Jehle-Reny
\newcommand{\calG}{\textbf{\textcolor{verde}{[C]}}}   % consolidação

% --- Notação canônica do projeto (PT-BR) ---
\newcommand{\TMS}{\text{TMS}}
% preferência denotada \succeq (já é o padrão amssymb; não MRS, não \succsim)
\DeclareMathOperator{\Esp}{\mathbb{E}}

% --- Decimais com vírgula: use \dc{0}{25} -> 0{,}25 ---
\newcommand{\dc}[2]{#1{,}#2}

% --- Caixa pedagógica (gabarito 5 passos) ---
\newtcolorbox{gabarito}{
  colback=insper!4, colframe=insper, boxrule=1pt, arc=3pt,
  left=6pt, right=6pt, top=6pt, bottom=6pt,
  breakable, before skip=8pt, after skip=8pt}

\newcommand{\passo}[1]{\smallskip\noindent\textbf{\textcolor{insper}{#1}}\quad}

% --- Caixa de ANDAIME DE LEITURA (explicação para não-economista) ---
% IMPORTANTE: esta caixa é material de ESTUDO EM CASA da versão-solução.
% NÃO consome orçamento de lousa: o professor escreve a álgebra (160 min
% preservados); o aluno usa a prosa azul para destravar cada passo sozinho.
% Cor azul-mundo (#1B3A5C) de propósito, para o leitor distinguir
% "andaime" (azul) de "comentário do professor" (vermelho gabarito).
\newtcolorbox{empalavras}[1][Em palavras (para quem não é economista)]{
  colback=mundo!5, colframe=mundo, boxrule=0.7pt, arc=2pt,
  left=6pt, right=6pt, top=5pt, bottom=5pt,
  fonttitle=\bfseries\small\sffamily, coltitle=white,
  title={#1}, breakable, before skip=6pt, after skip=6pt}

% --- Modo SOLUÇÃO (default) vs. CADERNO DO ALUNO ---
% Compilar este arquivo direto => \ifsolucao verdadeiro (com gabarito).
% Wrapper caderno-aluno-aula8.tex define \modoaluno => só enunciados.
\newif\ifsolucao
\ifdefined\modoaluno\solucaofalse\else\solucaotrue\fi

% Espaço de resolução para o caderno do aluno
\newcommand{\espacoaluno}[1]{%
  \par\smallskip\noindent\textit{\textcolor{gray}{Espaço para resolução.}}%
  \par\vspace{#1}\noindent\textcolor{insper!30}{\hrulefill}\par\medskip}

\title{\textbf{\textcolor{insper}{Microeconomia I — MPE 2026/32}}\\[4pt]
\large Aula 8 — Seleção Adversa e Risco Moral\\[2pt]
\normalsize \ifsolucao Três exercícios resolvidos para a lousa (revisão presencial)%
\else Caderno do aluno — três exercícios (enunciados para resolver)\fi}
\author{Insper — Mestrado Profissional em Economia}
\date{Calibre: Nicholson \& Snyder 12e Cap.~18 (piso) · Jehle-Reny 3e (teto)}

\begin{document}
\maketitle

\noindent\textbf{Convenções.} Documento em português do Brasil. Decimais com vírgula
($\dc{0}{25}$). Preferência denotada $\succeq$; taxa marginal de substituição $\TMS$.
Bernoulli do agente avesso $u(w)=\sqrt{w}$ (a raiz quadrada foi escolhida para que
\emph{todas} as contas fechem em inteiros na lousa). Os três exercícios são
\textbf{complementares} aos avaliativos da plataforma (não os duplicam): os números
foram recalibrados de propósito. \textbf{Calibre:} \calG{} consolidação,
\calY{} intermediário (Nicholson \& Snyder), \calR{} avançado (Jehle-Reny / qualifier).
Estes marcadores substituem os emojis verde/amarelo/vermelho do projeto, para compilar em pdflatex.

\ifsolucao
\smallskip
\noindent\textbf{\textcolor{mundo}{Como ler este material.}} A turma do MPE é heterogênea:
há advogados, engenheiros e gente de negócios que nunca viu microeconomia avançada. Por
isso a versão-solução traz, ao lado da álgebra, \textbf{caixas azuis ``Em palavras''}.
Elas são \emph{andaime de leitura para estudo em casa}: explicam cada passo em português
comum, definem o jargão na primeira aparição e traduzem cada resultado para a intuição.
Na lousa, o professor escreve apenas a álgebra (a caixa azul \emph{não} consome tempo de
quadro); você usa a prosa azul depois, sozinho, para destravar o que ficou opaco. As
\textbf{caixas vermelhas} (``Comentário pedagógico'') são a leitura do professor em 5 passos
— mantêm o registro técnico.
\fi

\medskip
\begin{center}
\begin{tabular}{@{}clcl@{}}
\toprule
\# & Tema (conceito-âncora da Aula 8) & Calibre & Lousa \\
\midrule
1 & Akerlof: \emph{hidden type}, unraveling, média condicional & \calY & $\sim$40 min \\
2 & Holmström: \emph{hidden action}, first/second-best, risco moral & \calR & $\sim$55 min \\
3 & Rothschild-Stiglitz: seguro, separador, single-crossing & \calY/\calR & $\sim$45 min \\
\bottomrule
\end{tabular}
\end{center}

\hrulefill

% =====================================================================
\section*{\textcolor{insper}{Exercício 1 — Akerlof: quando o mercado encolhe e quando colapsa}}
\noindent\textbf{Calibre \calY{} (N\&S §18.7).} \quad \emph{Lousa: $\sim$40 min.}
\par\smallskip
\noindent\textbf{Conceito que resume:} \emph{hidden type} (seleção adversa); preço de
equilíbrio competitivo igual à valoração esperada condicional ao conjunto de
vendedores ativos; o ``unraveling'' de Akerlof (1970); a fronteira entre
\emph{mercado parcial} e \emph{colapso total}.

\subsection*{Enunciado}
Há um continuum de carros usados. A qualidade $\theta$ é \textbf{informação privada do
vendedor} e se distribui $\theta \sim U[0,1]$. O dono de um carro de qualidade $\theta$
tem valor de reserva $\theta$ (não vende por menos). O comprador é competitivo, paga um
preço único $p$ e, por concorrência, esse preço iguala a \emph{valoração esperada} do
carro que ele consegue comprar (lucro esperado nulo).

\smallskip
\noindent\textbf{Parte A (aditiva).} O comprador valoriza cada carro em $\theta + k$, com
$k = \dc{0}{3}$ (um ``prêmio de uso'' fixo: revenda, prazer de dirigir etc.).

\smallskip
\noindent\textbf{Parte B (multiplicativa).} Agora o comprador valoriza $\beta\theta$ com
$\beta = \dc{1}{5}$. Compare o destino do mercado.

\ifsolucao

\begin{empalavras}[O enunciado em palavras de todo dia]
Imagine um site de carros usados. Cada carro tem uma \textbf{qualidade} verdadeira,
que chamamos de $\theta$ (a letra grega ``teta''). Só o \emph{dono} sabe o $\theta$ do
carro dele — o comprador não consegue olhar e descobrir; é \textbf{informação privada do
vendedor}. Dizer que $\theta \sim U[0,1]$ (``uniforme entre 0 e 1'') é só dizer: as
qualidades vão de 0 (sucata) a 1 (impecável), e qualquer valor no meio é igualmente
provável — como um número sorteado ao acaso entre 0 e 1.

\smallskip\noindent
O \textbf{valor de reserva} é o preço mínimo pelo qual o dono topa vender. Aqui o dono de
um carro de qualidade $\theta$ tem valor de reserva igual a $\theta$: quanto melhor o
carro, mais ele exige para se desfazer dele. Então, a um preço de mercado $p$, \emph{só}
os donos cujo carro vale $\theta \le p$ aceitam vender — os outros guardam o carro.

\smallskip\noindent
``Comprador competitivo'' + ``lucro esperado nulo'' é a \textbf{condição de livre entrada}:
como há muitos compradores (revendedores, particulares) disputando os carros, a concorrência
empurra o lucro a zero. Na prática isso significa: o preço pago acaba sendo \emph{exatamente}
quanto o comprador espera que o carro valha, em média, dado quem topa vender por aquele preço.
Nem mais (senão o comprador tem prejuízo esperado), nem menos (senão um rival oferece um
pouco mais e leva o negócio).
\end{empalavras}

\subsection*{Resolução passo a passo}
\begin{enumerate}[leftmargin=2em,itemsep=4pt]

\item \textbf{Benchmark simétrico (ferve em 1 linha).} Se o comprador observasse
$\theta$, haveria comércio sempre que a valoração supera a reserva. Na Parte A,
$\theta + k > \theta$ para todo $\theta$ (pois $k>0$): \textbf{todos} os carros
transacionam, ganho social $k$ por carro. Esse é o first-best.

\item \textbf{Quem está ativo a um preço $p$ (Parte A).} O vendedor só aceita se
$p \ge \theta$. Logo o conjunto ativo é $\{\theta : \theta \le p\}$ — exatamente os
\emph{piores} carros. Sob $U[0,1]$, condicionando em $\theta \le p$, a distribuição é
$U[0,p]$ e a qualidade média condicional é
\[
\Esp[\theta \mid \theta \le p] \;=\; \frac{p}{2}. \tag{1.1}
\]
\textbf{Aqui mora o veneno:} ao subir o preço, entram carros melhores, mas a média sobe
só pela metade da velocidade do preço.

\begin{empalavras}[Por que a média condicional é $p/2$?]
``Média condicional'' significa: a média da qualidade \emph{entre os carros que de fato
estão à venda} (não entre todos os carros do mundo). A condição é ``$\theta \le p$''.

\smallskip\noindent
Se a um preço $p$ só aparecem os carros com qualidade entre $0$ e $p$, e dentro desse
intervalo todas as qualidades são igualmente prováveis (continua uniforme, agora um
$U[0,p]$), então a média é simplesmente o \emph{ponto do meio} do intervalo: $\dfrac{0+p}{2}
= \dfrac{p}{2}$. Exemplo concreto: se $p = \dc{0}{6}$, os carros à venda têm qualidade
entre $0$ e $\dc{0}{6}$, e em média valem $\dc{0}{3}$.

\smallskip\noindent
O ``veneno'' citado ao lado: se eu, comprador, subo minha oferta de $\dc{0}{6}$ para
$\dc{0}{8}$, atraio carros um pouco melhores — mas a média deles sobe só de $\dc{0}{3}$
para $\dc{0}{4}$. O preço subiu $\dc{0}{2}$; a qualidade média subiu só $\dc{0}{1}$.
A média ``corre atrás'' do preço, mas com metade da velocidade. Esse descompasso é o
motor de tudo o que vem a seguir.
\end{empalavras}

\item \textbf{Equilíbrio (Parte A).} O comprador competitivo paga sua valoração esperada
condicional ao conjunto ativo:
\[
p \;=\; \Esp[\theta \mid \theta \le p] + k \;=\; \frac{p}{2} + k
\;\;\Longrightarrow\;\; \frac{p}{2} = k
\;\;\Longrightarrow\;\; \boxed{\,p^\ast = 2k = \dc{0}{6}\,}. \tag{1.2}
\]
Só transacionam os carros com $\theta \le \dc{0}{6}$ — \textbf{60\% do estoque}. Os 40\%
melhores ($\theta > \dc{0}{6}$) \emph{somem do mercado}. Qualidade média negociada:
$\dc{0}{3}$.

\begin{empalavras}[O que é ``ponto fixo'' e o que $p^\ast = \dc{0}{6}$ quer dizer]
Repare na estrutura circular: o preço depende de quem vende, mas quem vende depende do
preço. Para o mercado ``parar'' (entrar em equilíbrio), precisamos de um preço que seja
\emph{consistente consigo mesmo}: se o preço é $p$, então a média dos carros que aparecem é
$p/2$, então o comprador (que paga essa média mais o prêmio $k$) oferece $p/2 + k$ — e
esse valor tem de ser igual ao próprio $p$ de onde partimos. Um valor que ``volta nele
mesmo'' assim se chama \textbf{ponto fixo}. Resolver $p/2 + k = p$ acha esse ponto.

\smallskip\noindent
\textbf{Tradução do resultado:} $p^\ast = \dc{0}{6}$ não é só um número. Ele diz que
\emph{40\% dos melhores carros desaparecem do mercado} — e desaparecem \textbf{justamente
por serem bons}: seus donos exigem mais do que $\dc{0}{6}$, ninguém oferece tanto (porque o
comprador, em média, só encontra carros medianos), e o negócio não acontece. A assimetria
de informação \emph{expulsa a qualidade}. É a isto que Akerlof deu o nome de
\emph{unraveling} (algo como ``desfiamento'': o topo do mercado se desfaz).
\end{empalavras}

\item \textbf{Perda de bem-estar (Parte A).} Cada carro não-negociado destrói um ganho de
comércio igual a $k$. A massa não-negociada é $1 - p^\ast = \dc{0}{4}$, logo
\[
\text{perda} = k\,(1 - 2k) = \dc{0}{3}\cdot \dc{0}{4} = \dc{0}{12}. \tag{1.3}
\]
\textbf{Unraveling parcial:} o mercado não morre, mas encolhe — e some justamente o topo
da qualidade.

\begin{empalavras}[A ``perda de bem-estar'' em uma frase]
``Ganho de comércio'' é o valor que uma troca cria: o comprador valoriza o carro em
$\theta + k$, o dono só exige $\theta$, então cada negócio fechado gera $k = \dc{0}{3}$ de
valor novo para a sociedade (o prazer/uso extra que o comprador tira e o dono não tinha).
Quando um carro \emph{não} é negociado por causa da assimetria, esse $k$ simplesmente
\textbf{deixa de existir} — ninguém o captura. Como $\dc{0}{4}$ do estoque (os bons) ficam
parados, a conta da riqueza perdida é $\dc{0}{3} \times \dc{0}{4} = \dc{0}{12}$. É dinheiro
que evaporou só porque comprador e vendedor não conseguem confiar na qualidade.
\end{empalavras}

\item \textbf{Parte B (multiplicativa) — o colapso.} Agora a valoração é $\beta\theta$.
Conjunto ativo idêntico ($\theta \le p$), média condicional $p/2$. Equilíbrio:
\[
p = \beta\,\Esp[\theta\mid\theta\le p] = \beta\,\frac{p}{2}
\;\;\Longrightarrow\;\; p\Big(1 - \frac{\beta}{2}\Big) = 0. \tag{1.4}
\]
Como $\beta = \dc{1}{5} < 2$, o fator $1 - \beta/2 = \dc{0}{25} \ne 0$, então a
\textbf{única} solução é $p^\ast = 0$: \textbf{colapso total}. Nenhum carro é negociado.
A condição-limite é $\beta = 2$ (continuum degenerado); $\beta > 2$ resgata o mercado.

\begin{empalavras}[Por que ``vezes'' colapsa e ``mais'' não]
A diferença entre a Parte A e a Parte B é só \emph{como} o comprador valoriza o carro, e ela
muda tudo. Na Parte A (\textbf{aditiva}), o comprador paga ``qualidade $+\,k$'': há um bônus
$k = \dc{0}{3}$ \emph{fixo} que sobrevive mesmo que o carro seja ruim. Esse colchão fixo
mantém algum mercado vivo.

\smallskip\noindent
Na Parte B (\textbf{multiplicativa}), o comprador paga ``$\beta$ vezes a qualidade''. Se a
qualidade média cai, a valoração cai \emph{junto, proporcionalmente} — não há colchão. Faça
a conta no detalhe: para o mercado existir, o preço $p$ teria de igualar $\beta \cdot
(p/2)$. Isso só pode acontecer se $\beta/2 = 1$, ou seja $\beta = 2$. Como aqui $\beta =
\dc{1}{5}$ é \emph{menor} que $2$, cada vez que tento um preço positivo o comprador valoriza
o lote em \emph{menos} do que esse preço, então recua, o preço cai, somem mais carros bons,
a média piora, e a espiral só para em $p = 0$: \textbf{o mercado morre por inteiro}.

\smallskip\noindent
A lição: a seleção adversa nem sempre destrói o mercado (Parte A) — às vezes só corta o
topo. Quando ela destrói (Parte B) depende de a valoração do comprador ter ou não uma
parcela que não encolhe junto com a qualidade. ``$\beta > 1$'' (comprador valoriza mais que
o dono) \emph{não} basta; é preciso $\beta > 2$ por causa da divisão por dois da média.
\end{empalavras}

\end{enumerate}

\subsection*{Comentário pedagógico (5 passos)}
\begin{gabarito}
\passo{1. Ponto-chave.} O preço de equilíbrio resolve o ponto fixo
$p = \text{valoração}(\Esp[\theta\mid\theta\le p])$. A forma da valoração decide o
destino: \textbf{aditiva} ($\theta+k$) preserva um ganho fixo $k$ por carro $\Rightarrow$
mercado parcial; \textbf{multiplicativa} ($\beta\theta$) faz o ganho encolher junto com a
qualidade $\Rightarrow$ colapso quando $\beta<2$.

\passo{2. Derivação.} Tudo depende de $\Esp[\theta\mid\theta\le p]=p/2$ (média de um
uniforme truncado). Substituir e resolver o ponto fixo: $p/2+k=p$ dá $p^\ast=2k$;
$\beta p/2 = p$ dá $p(1-\beta/2)=0$. Uma conta de meia linha em cada caso — exatamente o
que se quer na lousa.

\passo{3. Armadilha.} (i) Achar que ``$\beta>1$ basta para o mercado funcionar'' — não:
na forma multiplicativa o que importa é $\beta$ \emph{vs.}~$2$, porque a média condicional
divide por dois. (ii) Confundir reserva do vendedor com preço de equilíbrio. (iii) Esquecer
que o conjunto ativo é \emph{truncado para baixo} ($\theta\le p$), o que é a essência da
seleção \emph{adversa}: quem fica é o pior.

\passo{4. Extensão.} O remédio é \emph{revelar tipo}. Reputação e garantia (Klein-Leffler
1981) atenuam o unraveling; \textbf{na Aula 9}, Spence (1973) introduz a sinalização — o
agente \emph{informado} (o bom vendedor) age para se separar do limão. Akerlof, Spence e
Stiglitz dividiram o Nobel de 2001 justamente por isso.

\passo{5. Peer-review (auto-crítica honesta).} A Parte A é a inovação pedagógica útil:
o aluno costuma sair da aula achando que ``Akerlof $=$ colapso''. A forma aditiva mostra
que \emph{nem sempre} colapsa — às vezes só encolhe pelo topo, o que é empiricamente o
caso típico (carros usados não desapareceram). Risco residual: um aluno apressado pode
aplicar a fórmula multiplicativa ($p^\ast=0$) na Parte A por hábito; o item~3 da resolução
deixa o contraste explícito. Não usei a forma multiplicativa com $\beta=\dc{1}{3}$ de
propósito — esse número já aparece nos avaliativos; aqui está $\beta=\dc{1}{5}$.
\end{gabarito}
\else
\espacoaluno{75mm}
\fi

\hrulefill

% =====================================================================
\section*{\textcolor{insper}{Exercício 2 — Holmström: o preço de não enxergar o esforço}}
\noindent\textbf{Calibre \calR{} (parte (a) é \calY{} de aquecimento).} \quad \emph{Lousa: $\sim$55 min.}
\par\smallskip
\noindent\textbf{Conceito que resume:} \emph{hidden action} (risco moral); principal-agente
de Holmström (1979); first-best (salário fixo, principal neutro absorve o risco) \emph{vs.}
second-best (IR + IC binding, $w_H > w_L$); \textbf{prêmio de risco moral}; e a
equivalência first-best $=$ second-best quando o agente é neutro ao risco.

\subsection*{Enunciado}
O Principal (P, \textbf{neutro} ao risco) contrata um Agente (A, \textbf{avesso},
Bernoulli $u(w)=\sqrt{w}$). O esforço $a \in \{a_L, a_H\}$ é privado do agente, com custos
$c(a_L)=0$ e $c(a_H)=2$. O produto observável é $y \in \{0,\,80\}$. A probabilidade de
$y=80$ é $\pi_H = \dc{0}{8}$ sob esforço alto e $\pi_L = \dc{0}{4}$ sob esforço baixo. A
utilidade de reserva do agente é $\bar U = 3$. O contrato é um salário $w(y)$ condicionado
ao produto. \textbf{O Principal deseja induzir $a_H$} (verifique que é eficiente fazê-lo).

\ifsolucao

\begin{empalavras}[O enunciado em palavras de todo dia]
Pense num \textbf{patrão} (o ``Principal'', P) que contrata um \textbf{funcionário} (o
``Agente'', A) para uma tarefa. O funcionário pode se esforçar muito ($a_H$, esforço alto)
ou pouco ($a_L$, esforço baixo). Esforçar-se \emph{custa} para ele: $c(a_H) = 2$ unidades de
incômodo (cansaço, abrir mão do tempo livre), enquanto não se esforçar custa $0$. O problema:
o patrão \textbf{não vê} se o funcionário se esforçou — é \emph{ação oculta} (``hidden
action''). Ele só vê o \emph{resultado} $y$: ou o projeto dá certo ($y = 80$) ou dá errado
($y = 0$).

\smallskip\noindent
Esforço não garante sucesso, só aumenta a \emph{chance}: com esforço alto, o projeto dá
certo com probabilidade $\pi_H = \dc{0}{8}$; com esforço baixo, só $\pi_L = \dc{0}{4}$.

\smallskip\noindent
\textbf{``Neutro'' vs.\ ``avesso'' ao risco.} O patrão é \emph{neutro ao risco}: só liga
para o valor médio do dinheiro, tanto faz se vem garantido ou no susto. O funcionário é
\emph{avesso ao risco}: ele prefere um salário garantido a uma aposta de mesmo valor
médio — como quase todos nós, que preferimos um salário fixo a ganhar ``às vezes muito, às
vezes nada''. Modelamos isso pela \textbf{função de utilidade de Bernoulli} $u(w) = \sqrt{w}$
(``utilidade sobre a riqueza''): a raiz quadrada cresce \emph{cada vez mais devagar}, então
cada real a mais traz menos satisfação que o anterior — é exatamente o que significa ter
medo de risco (ver a próxima caixa). $\bar U = 3$ é a \textbf{utilidade de reserva}: o
quanto de satisfação o funcionário consegue em outro emprego; se o contrato não der pelo
menos isso, ele vai embora.
\end{empalavras}

\subsection*{Resolução passo a passo}
\begin{enumerate}[leftmargin=2em,itemsep=4pt]

\item \textbf{(a) First-best: ação observável.} Se P enxerga $a$, ele exige $a_H$ por
contrato e paga um \textbf{salário fixo} $w^\ast$ (independente de $y$) que satura a
participação. A restrição IR binding é
\[
\sqrt{w^\ast} - c(a_H) = \bar U
\;\;\Longrightarrow\;\; \sqrt{w^\ast} = \bar U + c(a_H) = 3 + 2 = 5
\;\;\Longrightarrow\;\; \boxed{\,w^\ast = 25\,}. \tag{2.1}
\]
Salário fixo: \textbf{o principal neutro carrega todo o risco}, e o agente avesso vive sem
risco. É Pareto-eficiente.

\begin{empalavras}[Por que a raiz quadrada modela ``medo de risco'', e o que é IR]
``First-best'' (``melhor caso'') é o mundo de sonho em que o patrão \emph{vê} o esforço.
Nesse mundo ele não precisa de truque: escreve no contrato ``faça esforço alto'' e paga um
\textbf{salário fixo} $w^\ast$, igual dê certo ou errado. Como o funcionário odeia risco e o
patrão não liga, o arranjo ótimo é o patrão \emph{absorver toda a incerteza} e dar ao
funcionário um pagamento garantido. É a divisão de risco perfeita.

\smallskip\noindent
\textbf{Por que $\sqrt{w}$ é ``avesso ao risco''?} Compare duas situações de mesmo valor
médio. (i) Receber $25$ garantidos: satisfação $\sqrt{25} = 5$. (ii) Apostar: metade das
vezes $0$, metade das vezes $50$ (média também $25$): satisfação esperada $\tfrac12\sqrt{0}
+ \tfrac12\sqrt{50} \approx \tfrac12(0) + \tfrac12(7{,}07) \approx 3{,}5$. O valor médio é o
mesmo ($25$), mas a aposta dá \emph{menos} satisfação ($3{,}5 < 5$). É a curvatura da raiz
(``cada real extra satisfaz menos'') que faz a pessoa fugir do risco. Guarde isto: é a razão
de o salário fixo ser tão valioso para o funcionário.

\smallskip\noindent
\textbf{A restrição IR (``participação'').} ``IR'' vem de \emph{individual rationality}:
o funcionário só aceita o emprego se ele render \emph{pelo menos} a satisfação que teria
fora ($\bar U = 3$). A equação $\sqrt{w^\ast} - c(a_H) = \bar U$ diz: a satisfação do
salário, \emph{descontado} o incômodo do esforço, tem de empatar com o emprego de fora.
Dizemos que a IR está \textbf{binding} (``apertada'', ``ativa'') quando vale com igualdade
($=$, não $>$): o patrão não paga \emph{nenhum} centavo a mais do que o necessário para o
funcionário não ir embora — afinal, ele quer gastar o mínimo. Daí $\sqrt{w^\ast} = 3 + 2 =
5$, logo $w^\ast = 25$.
\end{empalavras}

\item \textbf{Checagem de que $a_H$ é eficiente.} Lucro esperado do P sob first-best:
$\Esp[y\mid a_H] - w^\ast = \dc{0}{8}\cdot 80 - 25 = 64 - 25 = 39$. Sob $a_L$, o salário
fixo seria $\sqrt{w} = \bar U = 3 \Rightarrow w = 9$, e o lucro
$\dc{0}{4}\cdot 80 - 9 = 32 - 9 = 23$. Como $39 > 23$, induzir $a_H$ é eficiente.
\textbf{Guarde o número $39$.}

\item \textbf{(b) Por que salário fixo quebra sob ação oculta.} Se P não vê $a$ e paga
fixo, o agente recebe $\sqrt{w}$ faça o que fizer; como $c(a_H) > c(a_L)=0$, ele
\textbf{desvia} para $a_L$ (mesmo salário, menos custo). A IC para $a_H$ é violada. Para
induzir esforço, o salário \emph{tem} de depender de $y$.

\item \textbf{(c) Second-best: o programa.} P minimiza o salário esperado sujeito a
participar (IR) e a não desviar (IC):
\[
\min_{w_H,\,w_L}\;\; \pi_H w_H + (1-\pi_H) w_L
\]
\[
\text{(IR)}\quad \pi_H\sqrt{w_H} + (1-\pi_H)\sqrt{w_L} - c(a_H) \ge \bar U,
\]
\[
\text{(IC)}\quad (\pi_H - \pi_L)\big[\sqrt{w_H} - \sqrt{w_L}\big] \ge c(a_H) - c(a_L).
\]
Com agente avesso e $a_H$ desejado, \textbf{ambas} as restrições são binding (resultado
padrão; o gabarito justifica). Substituindo os números:

\begin{empalavras}[O que são IR e IC, e por que as duas ficam ``apertadas'']
Agora estamos no \textbf{second-best} (``segundo melhor''): o mundo real, em que o patrão
\emph{não} vê o esforço. Salário fixo não funciona mais — se o pagamento é o mesmo dê no que
der, o funcionário escolhe não se esforçar (mesmo dinheiro, menos cansaço). Para induzir
esforço, o patrão precisa \textbf{premiar o bom resultado}: pagar mais quando $y = 80$
($w_H$) do que quando $y = 0$ ($w_L$). É o \emph{bônus por desempenho}.

\smallskip\noindent
O patrão monta um problema de minimização: ``qual o menor salário esperado que eu consigo
pagar, respeitando duas condições?'' As duas condições são:
\begin{itemize}[leftmargin=1.3em,itemsep=1pt,topsep=2pt]
\item \textbf{IR (participação):} o funcionário tem de aceitar o emprego — sua satisfação
esperada, menos o custo do esforço, $\ge \bar U$. Sem isso, ele vai embora.
\item \textbf{IC (incentivo, de \emph{incentive compatibility}):} dado o contrato, o
funcionário tem de \emph{preferir por conta própria} se esforçar a vacilar. A IC compara a
satisfação esperada esforçando-se com a de não se esforçar; o esforço só ``compensa'' se o
prêmio extra por desempenho cobrir o custo do esforço.
\end{itemize}

\smallskip\noindent
\textbf{Por que as duas ficam binding (com igualdade)?} Intuição de patrão sovina: ele
nunca dá nada de graça. (i) Se a IR sobrasse (funcionário satisfeito \emph{além} do
necessário), o patrão cortaria salário até encostar em $\bar U$ — logo IR vira igualdade.
(ii) Se a IC sobrasse (funcionário esforçando-se com folga de incentivo), o patrão
\emph{encolheria o bônus} ($w_H - w_L$), o que reduz o risco imposto ao funcionário avesso e
\emph{barateia} a folha — até a IC encostar na igualdade. Como mexer em cada uma sempre
reduz custo até o limite, no ótimo \textbf{ambas valem com $=$}. É por isso que trocamos os
$\ge$ por $=$ nas contas a seguir.
\end{empalavras}

\[
\text{(IR)}\quad \dc{0}{8}\sqrt{w_H} + \dc{0}{2}\sqrt{w_L} = 3 + 2 = 5, \tag{2.2}
\]
\[
\text{(IC)}\quad (\dc{0}{8} - \dc{0}{4})\big[\sqrt{w_H} - \sqrt{w_L}\big] = 2
\;\Longrightarrow\; \dc{0}{4}\big[\sqrt{w_H}-\sqrt{w_L}\big] = 2
\;\Longrightarrow\; \sqrt{w_H} - \sqrt{w_L} = 5. \tag{2.3}
\]

\item \textbf{(c) Resolver o sistema $2\times2$ em $(\sqrt{w_H},\sqrt{w_L})$.} De (2.3),
$\sqrt{w_H} = 5 + \sqrt{w_L}$. Substituindo em (2.2):
\[
\dc{0}{8}\,(5 + \sqrt{w_L}) + \dc{0}{2}\,\sqrt{w_L} = 5
\;\Longrightarrow\; 4 + \sqrt{w_L} = 5
\;\Longrightarrow\; \sqrt{w_L} = 1.
\]
Logo $\sqrt{w_L}=1 \Rightarrow w_L = 1$, e $\sqrt{w_H} = 6 \Rightarrow w_H = 36$.
\[
\boxed{\,w_L = 1, \qquad w_H = 36\,}. \tag{2.4}
\]

\item \textbf{(d) Prêmio de risco moral.} O custo esperado da folha sob second-best:
\[
\Esp[w] = \pi_H w_H + (1-\pi_H) w_L = \dc{0}{8}\cdot 36 + \dc{0}{2}\cdot 1
= \dc{28}{8} + \dc{0}{2} = 29.
\]
Compare com o first-best ($w^\ast = 25$):
\[
\boxed{\;\text{prêmio de risco moral} = \Esp[w]_{\text{SB}} - w^\ast_{\text{FB}}
= 29 - 25 = 4.\;} \tag{2.5}
\]
São 4 unidades a mais de folha — a ``sombra'' de não enxergar o esforço. O lucro do P cai
de $39$ (first-best) para $64 - 29 = 35$ (second-best); a queda de $4$ é exatamente o
prêmio.

\begin{empalavras}[O que significa ``prêmio de risco moral'' = 4]
Compare os dois mundos. No \emph{first-best} o patrão pagava $25$ garantidos. No mundo real
ele precisa pagar, \emph{em média}, $29$ — porque agora dá um bônus arriscado ($w_H = 36$ se
dá certo, só $w_L = 1$ se dá errado), e o funcionário avesso \emph{cobra mais na média} para
topar carregar esse risco. A diferença, $29 - 25 = 4$, é o \textbf{prêmio de risco moral}:
o sobrepreço que o patrão paga só porque não consegue enxergar o esforço e teve de comprar
incentivo com risco.

\smallskip\noindent
Note a simetria: o lucro do patrão também cai exatamente $4$ (de $39$ para $35$). Faz
sentido — o valor que o projeto gera ($64$ em média) não mudou; o que mudou foi que $4$
unidades a mais saíram do bolso do patrão na forma de folha. \textbf{Atenção a uma confusão
clássica:} o funcionário \emph{não} fica $4$ mais rico. A satisfação dele continua em
$\bar U = 3$, igual ao first-best — a IR aperta nos dois mundos. Os $4$ extras de salário
esperado são \emph{inteiramente devorados} pelo desprazer de carregar o bônus arriscado:
não são presente a ninguém. Por isso esse $4$ é um \emph{custo da assimetria de informação}
— riqueza que \textbf{evapora} (peso morto): quem perde os $4$ é o patrão, e quem ganha
é\ldots\ ninguém.
Em CEO, é o desconto que acionistas pagam ao dar \emph{stock options} arriscadas em vez de
salário fixo.
\end{empalavras}

\item \textbf{(e) Equivalência sob neutralidade do agente.} Se o agente também fosse
neutro ($u(w)=w$), o P poderia \textbf{vender o resultado} ao agente: cobra uma taxa fixa
$F$ e deixa o agente como \emph{residual claimant} de $y$. O agente neutro escolhe $a_H$
por interesse próprio (maximiza $\Esp[y\mid a]-c(a)$), sem custo de risco. Resultado:
\[
\text{second-best} = \text{first-best}, \qquad \text{prêmio} = 0.
\]
O prêmio de risco moral \textbf{nasce especificamente da aversão a risco do agente} — é o
preço de ``quebrar o seguro'' do agente para criar incentivo.

\begin{empalavras}[``Vender o resultado'' e o que é \emph{residual claimant}]
Aqui o exercício pergunta: e se o funcionário \emph{também} não tivesse medo de risco
(fosse neutro, $u(w) = w$)? Então o patrão pode usar um truque elegante, o contrato de
\emph{franquia} (``franchise''): em vez de assalariar, ele \textbf{vende a operação} ao
funcionário por uma taxa fixa $F$ e deixa que o funcionário fique com \emph{tudo} o que
sobrar do resultado. Quem fica com o que sobra chama-se \textbf{\emph{residual claimant}}
(``reclamante residual'', o dono do lucro líquido). O exemplo de manual é o
\textbf{arrendamento de terra por renda fixa}: o arrendatário paga um \emph{aluguel fixo}
(o nosso $F$) e fica com \emph{toda} a colheita — e, sem medo de risco, planta o ponto
eficiente por conta própria. O contraste fecha a ideia: a \emph{meiação} (dividir a colheita
com o dono da terra) só aparece quando o arrendatário \emph{teme} risco, trocando incentivo
por seguro — exatamente como nos itens anteriores. (Franquias reais, tipo McDonald's, cobram
um \emph{percentual} sobre as vendas justamente por isso: o franqueado de carne e osso
\emph{não} é neutro ao risco, então o contrato puro de taxa fixa não é o que se observa.)

\smallskip\noindent
Como agora o próprio bolso do funcionário absorve o resultado, ele \emph{se esforça por
interesse próprio}, sem que ninguém precise vigiá-lo — a IC se resolve sozinha. E como ele
não tem medo de risco, carregar a incerteza não lhe custa nada. Resultado: nenhum prêmio,
second-best $=$ first-best. \textbf{Moral da história:} o prêmio de risco moral dos itens
anteriores não nasce da assimetria \emph{em si}, e sim da combinação assimetria $+$ um
agente que \emph{teme risco}. Tire o medo de risco e o problema some.
\end{empalavras}

\end{enumerate}

\subsection*{Comentário pedagógico (5 passos)}
\begin{gabarito}
\passo{1. Ponto-chave.} Risco moral força o \emph{trade-off seguro vs.~incentivo}. O
first-best dá seguro pleno ($w$ fixo). Para induzir esforço sem observá-lo, o P precisa
\emph{premiar o produto alto} ($w_H>w_L$), e isso reintroduz risco no agente avesso, que
exige compensação na média. A diferença $\Esp[w]_{\text{SB}} - w_{\text{FB}}$ é o prêmio
de risco moral.

\passo{2. Derivação.} O sistema $2\times2$ em $(\sqrt{w_H},\sqrt{w_L})$ é o coração: IC
binding fixa a \emph{diferença} ($\sqrt{w_H}-\sqrt{w_L}=5$); IR binding fixa a
\emph{média ponderada} ($\dc{0}{8}\sqrt{w_H}+\dc{0}{2}\sqrt{w_L}=5$). Substituir uma na
outra resolve em duas linhas: $\sqrt{w_L}=1$, $\sqrt{w_H}=6$. Os parâmetros foram
calibrados para fechar em inteiros ($w_L=1$, $w_H=36$, $\Esp[w]=29$, prêmio $=4$).

\passo{3. Armadilha.} (i) Esquecer o coeficiente $(\pi_H-\pi_L)=\dc{0}{4}$ na IC e escrever
$\sqrt{w_H}-\sqrt{w_L}=2$ em vez de $5$ — erro fatal e comum. (ii) Esquecer de somar
$c(a_H)$ ao $\bar U$ na IR. (iii) Linearizar a Bernoulli (tratar $w$ no lugar de $\sqrt{w}$).
(iv) Resolver $\sqrt{w_L}=\sqrt{w_H}-5$ e aceitar raiz negativa. Os quatro erros mudam a
resposta — vale escrevê-los no quadro como ``não faça isto''.

\passo{4. Extensão.} A \emph{fórmula informativa} de Holmström (1979): o salário ótimo
depende de $y$ apenas via o quanto $y$ \emph{informa} sobre $a$ (razão de
verossimilhança). Aplicações: stock options para CEO (acionistas diversificados $=$ P
neutro; CEO $=$ A avesso); \emph{Bolsa Família} condicionada (frequência escolar
$\ge 85\%$ é o ``$y$ verificável'' que implementa a IC). \textbf{Aula 9} vira o problema do
avesso: lá é o agente \emph{informado} que age (sinalização).

\passo{5. Peer-review (auto-crítica honesta).} Os números são distintos dos avaliativos
(lá: $w_H=121$, $w_L=1$, prêmio $=21$); aqui o prêmio é só $4$ — \emph{menor de propósito},
para que a aritmética flua e o aluno enxergue o mecanismo sem se afogar em contas. Escolha
deliberada: aqui $a_H$ é eficiente \emph{também} no first-best (lucro $39>23$), o que torna
o second-best uma resposta a um problema ``de verdade'' (P quer mesmo o esforço alto) — nos
avaliativos o $a_L$ era first-best, o que confunde alguns alunos. Risco residual: a parte
(e) (neutralidade) é a peça destacável se o tempo apertar; ela é conceitual, não algébrica,
então cortá-la não fere a derivação central — mas perde o ``aha'' de que o problema é a
aversão, não a assimetria \emph{per se}.
\end{gabarito}
\else
\espacoaluno{85mm}
\fi

\hrulefill

% =====================================================================
\section*{\textcolor{insper}{Exercício 3 — Rothschild-Stiglitz: o bom risco paga a conta}}
\noindent\textbf{Calibre \calY{} com cauda \calR{} no fechamento algébrico.} \quad \emph{Lousa: $\sim$45 min.}
\par\smallskip
\noindent\textbf{Conceito que resume:} \emph{hidden type} em seguro; equilíbrio separador
de Rothschild-Stiglitz (1976); \textbf{single-crossing}; o tipo alto risco ($H$) fica com
cobertura completa, o tipo baixo risco ($L$) é \emph{distorcido para baixo} até que a IC do
$H$ fique binding; menção à inexistência sob $\lambda$ baixo.

\subsection*{Enunciado}
Mercado de seguro competitivo. Riqueza inicial $W = 17$; sinistro de tamanho $L = 16$
(quase tudo: sem seguro, o azarado fica com $W-L=1$). Há dois tipos: alto risco $H$ com
$\pi^H = \dc{0}{5}$ e baixo risco $L$ com $\pi^L = \dc{0}{25}$. O agente conhece o próprio
tipo; a seguradora não. Bernoulli $u(W)=\sqrt{W}$ (avesso). Um contrato é um par
$(p,c)$: \emph{prêmio} $p$ (paga sempre) e \emph{cobertura} $c\in[0,L]$ (recebe se há
sinistro). Por concorrência (livre entrada), cada contrato aceito dá lucro esperado nulo:
$p = \pi^\theta c$ no tipo que o compra.
\emph{(Números calibrados para fechamento exibível: $W=17,\,L=16$ é uma escolha
pedagógica deliberada — ver passo 5.)}

\ifsolucao

\begin{empalavras}[O enunciado em palavras de todo dia]
Pense em \textbf{seguro de carro} (ou plano de saúde). Você tem uma riqueza inicial $W =
17$. Existe o risco de um \textbf{sinistro} (bater o carro, adoecer) que custa $L = 16$ —
quase tudo: sem seguro, se der azar, você fica com só $W - L = 1$.

\smallskip\noindent
Há dois tipos de cliente. O \textbf{alto risco} ($H$) bate o carro com probabilidade
$\pi^H = \dc{0}{5}$ (meio a meio); o \textbf{baixo risco} ($L$) só com $\pi^L = \dc{0}{25}$
(motorista cuidadoso). \emph{Cada cliente sabe que tipo é}, mas a \textbf{seguradora não}
consegue distinguir — é informação privada (de novo, \emph{hidden type}). Cuidado com a
notação: a letra $L$ aparece em dois papéis — \emph{tamanho do sinistro} ($L = 16$) e
\emph{tipo baixo risco}. O contexto sempre deixa claro qual.

\smallskip\noindent
Um \textbf{contrato} é um par $(p, c)$: você paga um \textbf{prêmio} $p$ \emph{sempre} (deu
sinistro ou não — é a mensalidade do seguro) e recebe uma \textbf{cobertura} $c$ \emph{se}
houver sinistro. ``Lucro esperado nulo'' / \textbf{livre entrada} é a mesma ideia do
Exercício~1: concorrência entre seguradoras empurra o lucro a zero, então o prêmio cobre
exatamente o custo esperado, $p = \pi^\theta c$ (probabilidade de pagar $\times$ quanto paga).
O cliente é \emph{avesso ao risco} ($u(W) = \sqrt{W}$, mesma raiz quadrada do Exercício~2),
então quer suavizar: prefere consumo parecido nos dois cenários (com ou sem sinistro).
\end{empalavras}

\subsection*{Resolução passo a passo}
\begin{enumerate}[leftmargin=2em,itemsep=4pt]

\item \textbf{First-best simétrico.} Se a seguradora observasse o tipo, ofereceria a cada
um cobertura completa ($c=L$) ao prêmio justo $p^\theta = \pi^\theta L$. Razão: o agente
avesso quer suavizar consumo entre os estados, e a seguradora neutra absorve o risco a
custo zero esperado. Aqui:
\[
p^H = \pi^H L = \dc{0}{5}\cdot 16 = 8, \qquad
p^L = \pi^L L = \dc{0}{25}\cdot 16 = 4. \tag{3.1}
\]
Cada tipo fica com consumo \emph{certo}: $H$ tem $W-p^H = 9$; $L$ tem $W-p^L = 13$.

\begin{empalavras}[Por que ``cobertura completa'' e como sai o consumo certo]
Cobertura completa significa $c = L$: se der sinistro, o seguro repõe \emph{tudo} o que você
perdeu. Aí seu consumo fica \textbf{igual} nos dois cenários — você paga o prêmio $p$
sempre, e o sinistro é totalmente compensado. Por isso o consumo vira ``certo'' (sem
sobressalto): $W - p$ aconteça o que acontecer. Para o $H$, $17 - 8 = 9$; para o $L$, $17 -
4 = 13$. Como o cliente \emph{teme risco}, esse seguro pleno é o ideal para ele, e a
seguradora neutra topa porque, na média, sai no zero a zero. Esse é o \emph{first-best}: o
mundo em que a seguradora enxerga o tipo e cobra de cada um o prêmio justo.
\end{empalavras}

\item \textbf{Single-crossing (a chave geométrica).} Sob aversão e $\pi^H>\pi^L$, as
curvas de indiferença dos dois tipos no plano (prêmio $\times$ cobertura) se cruzam
\emph{exatamente uma vez}: na margem, o tipo $H$ (mais provável de sinistrar) está disposto
a pagar \emph{mais} por uma unidade extra de cobertura. É isso que permite \emph{separar}
os tipos com um cardápio de contratos.

\begin{empalavras}[``Single-crossing'' e ``separar tipos'' sem economês]
A seguradora não sabe quem é quem, então monta um \textbf{cardápio} de contratos e deixa
cada cliente \emph{escolher} — torcendo para que cada tipo escolha o contrato ``feito para
ele''. Isso se chama \textbf{separar} os tipos (cada um se revela pela escolha). Mas só
funciona se os dois tipos tiverem \emph{gostos diferentes o bastante} para não quererem o
mesmo contrato.

\smallskip\noindent
\textbf{Single-crossing} (``cruzamento único'') é a propriedade que garante isso. Pense em
quanto cada tipo valoriza ``um pedacinho a mais de cobertura''. O alto risco ($H$) bate o
carro com mais frequência, então uma cobertura extra é mais útil para ele — ele paga
\emph{mais caro} por ela do que o baixo risco. Tecnicamente: as curvas de indiferença (linhas
de ``mesma satisfação'') dos dois tipos, desenhadas no plano prêmio~$\times$~cobertura, se
cruzam \textbf{exatamente uma vez}. Esse cruzamento único é o que permite desenhar um
contrato que o $L$ aceita mas o $H$ recusa — a base de toda a separação. Sem single-crossing,
os gostos se embaralhariam e não daria para separar.
\end{empalavras}

\item \textbf{O separador R-S — estrutura.} O único candidato a equilíbrio competitivo é:
\begin{itemize}[leftmargin=1.4em,itemsep=2pt]
\item \textbf{Tipo $H$:} cobertura completa ao prêmio justo, $(p^H_\ast, c^H_\ast) = (8,16)$.
Sem distorção: o $H$ já não tem incentivo a imitar o $L$ (o $L$ tem cobertura menor).
\item \textbf{Tipo $L$:} cobertura \emph{parcial} $c^L_\ast < L$, sobre sua reta de prêmio
justo $p^L_\ast = \pi^L c^L_\ast = \dc{0}{25}\,c^L_\ast$, escolhida no ponto em que a
\textbf{IC do $H$ fica binding} — o $H$ é exatamente indiferente entre seu contrato e o do $L$.
\end{itemize}

\item \textbf{Montar a IC do $H$ binding.} O $H$ no seu contrato tem consumo certo
$W-p^H = 9$, logo utilidade $\sqrt{9}=3$. Se imitasse o contrato do $L$, enfrentaria risco:
\[
\underbrace{\sqrt{W - p^H}}_{\text{contrato do }H}
\;=\;
(1-\pi^H)\sqrt{\,W - p^L_\ast\,} + \pi^H \sqrt{\,W - p^L_\ast - L + c^L_\ast\,}.
\tag{3.2}
\]
Com $\pi^H=\dc{0}{5}$ e $p^L_\ast = \dc{0}{25}\,c^L_\ast$, e abreviando
\[
A \equiv W - p^L_\ast = 17 - \dc{0}{25}c^L_\ast, \qquad
B \equiv W - p^L_\ast - L + c^L_\ast = 1 + \dc{0}{75}c^L_\ast,
\]
a equação (3.2) vira (lado esquerdo $=3$):
\[
3 = \tfrac12\big(\sqrt{A} + \sqrt{B}\big)
\;\;\Longrightarrow\;\; \sqrt{A} + \sqrt{B} = 6. \tag{3.3}
\]

\begin{empalavras}[O que a equação (3.2) está dizendo de verdade]
Tradução da equação central: estamos \textbf{prendendo} o contrato do $L$ no maior tamanho
que ainda \emph{não tente} o $H$ a imitá-lo. O lado esquerdo, $\sqrt{W - p^H} = \sqrt{9} =
3$, é a satisfação do $H$ \emph{no contrato dele} (cobertura completa, consumo certo $9$).
O lado direito é a satisfação que o $H$ teria \emph{se mentisse} e pegasse o contrato barato
do $L$: aí ele ficaria \emph{exposto ao risco} — com probabilidade $1 - \pi^H$ não há
sinistro (consumo $A$) e com $\pi^H$ há sinistro (consumo $B$), daí a média ponderada das
duas raízes. Igualar os dois lados é dizer ``o $H$ é \emph{exatamente indiferente} entre ser
honesto e mentir''. Esse é o ponto em que a \textbf{IC do $H$ fica binding}: encolhemos a
cobertura do $L$ até bem ali, nem um milímetro a mais (senão o $H$ mentiria e a separação
quebra), nem a menos (senão estaríamos punindo o $L$ sem necessidade). $A$ e $B$ são só
apelidos para o consumo do $L$ nos estados ``sem sinistro'' e ``com sinistro''.
\end{empalavras}

\item \textbf{Fechamento algébrico (giz-limpo).} Note que
$A - B = (17 - \dc{0}{25}c^L_\ast) - (1 + \dc{0}{75}c^L_\ast) = 16 - c^L_\ast = L - c^L_\ast$.
Como $A-B = (\sqrt{A}-\sqrt{B})(\sqrt{A}+\sqrt{B})$ e $\sqrt{A}+\sqrt{B}=6$:
\[
\sqrt{A} - \sqrt{B} = \frac{A-B}{6} = \frac{16 - c^L_\ast}{6}. \tag{3.4}
\]
Somando (3.3) e (3.4): $2\sqrt{A} = 6 + \dfrac{16-c^L_\ast}{6}$. Em vez de elevar ao
quadrado o tempo todo, use o atalho: de (3.3) e (3.4), $\sqrt{A}\sqrt{B} =
\frac{1}{4}\big[(\sqrt{A}+\sqrt{B})^2 - (\sqrt{A}-\sqrt{B})^2\big]$, mas o caminho mais
curto é substituir um par-teste. \textbf{Tente $c^L_\ast = 4$:}
\[
A = 17 - \dc{0}{25}\cdot 4 = 16, \quad B = 1 + \dc{0}{75}\cdot 4 = 4,
\quad \sqrt{A}+\sqrt{B} = 4 + 2 = 6 \;\checkmark
\]
Bate (3.3) exatamente. Logo
\[
\boxed{\,c^L_\ast = 4, \qquad p^L_\ast = \dc{0}{25}\cdot 4 = 1\,}. \tag{3.5}
\]
\emph{(Quem preferir o caminho fechado: (3.3) e (3.4) dão a quadrática
$(c^L_\ast)^2 - 4W\,c^L_\ast + L^2 = 0$, isto é $(c^L_\ast)^2 - 68\,c^L_\ast + 256 = 0$,
discriminante $68^2 - 4\cdot256 = 4624 - 1024 = 3600 = 60^2$, raízes
$c^L_\ast = (68\pm 60)/2 \in \{4,\,64\}$; só $4$ satisfaz $c\le L=16$.)}

\begin{empalavras}[O truque do $\sqrt{A}-\sqrt{B}$ e por que dá para ``chutar'' $c = 4$]
A equação (3.3) tem \emph{duas} raízes quadradas somadas — incômodo de resolver na marra.
O atalho usa uma identidade da escola: $(\sqrt{A}-\sqrt{B})(\sqrt{A}+\sqrt{B}) = A - B$
(é só $a^2 - b^2 = (a-b)(a+b)$ com $a = \sqrt{A}$, $b = \sqrt{B}$). Já sabemos que a soma é
$6$ (de 3.3) e que a diferença $A - B = 16 - c$ (conta direta). Então a \emph{diferença} das
raízes é $\dfrac{A-B}{6}$. Com soma e diferença das raízes em mãos, descobrir cada raiz é
trivial. Esse malabarismo evita elevar ao quadrado duas vezes (que geraria termos cruzados
feios).

\smallskip\noindent
\textbf{E o ``chute'' $c = 4$?} Não é mágica nem sorte: os números $W = 17$, $L = 16$ foram
escolhidos pelo professor \emph{de propósito} para que a resposta caísse num inteiro redondo
e a verificação fosse instantânea. Substituir $c = 4$ dá $A = 16$ e $B = 4$, com $\sqrt{16}
+ \sqrt{4} = 4 + 2 = 6$ — bate na hora. Para o aluno cético que não confia em chute, o
parágrafo acima mostra o caminho ``oficial'': montar a equação do segundo grau $c^2 - 68c +
256 = 0$ e resolvê-la pela fórmula de Bhaskara. As duas rotas chegam ao mesmo $c = 4$ (a
outra raiz, $64$, é descartada porque a cobertura não pode passar do tamanho do sinistro,
$c \le L = 16$).
\end{empalavras}

\item \textbf{Leitura do resultado.} O tipo $L$ recebe $c^L_\ast = 4$ de $16$ possíveis:
\textbf{só 25\% de cobertura}. Ele carrega risco residual (consumo $16$ no estado bom,
$4$ no ruim, $\sqrt{16}\ne\sqrt{4}$). \textbf{O bom risco paga o preço da assimetria} —
aceita seguro pior do que mereceria para que o alto risco não o imite. Distorção de
second-best.

\item \textbf{Inexistência (só a intuição, sem formalizar).} Se a proporção $\lambda$ de
tipos $H$ é \emph{baixa}, a seguradora poderia oferecer um \emph{pooling} com cobertura
quase completa a prêmio médio (próximo de $\pi^L$), que Pareto-domina o separador. Mas
pooling \textbf{não é Nash}: uma rival desvia com um contrato de cobertura parcial mais
barato que atrai só os $L$, e o pool desaba. Nem separador nem pooling sobrevivem —
\textbf{R-S podem não ter equilíbrio competitivo}. (Não formalizamos: exige jogo contínuo.)

\begin{empalavras}[O resultado em uma frase, e por que o mercado pode não fechar]
\textbf{O paradoxo central:} quem ``paga a conta'' da assimetria é o \emph{bom} motorista
($L$), não o ruim. O motorista cuidadoso é obrigado a aceitar um seguro \emph{pior do que
mereceria} — só $4$ de $16$ de cobertura, $25\%$ — porque, se a seguradora lhe oferecesse a
cobertura cheia e barata que ele de fato merece, o motorista \emph{ruim} ($H$) correria para
comprá-la também (mentindo sobre seu tipo). Para evitar essa imitação, a seguradora ``capa''
o seguro do bom motorista. O alto risco, esse, sai ileso: cobertura completa, prêmio justo.
A intuição moralista (``o arriscado deveria levar menos seguro'') está \emph{invertida} —
é o oposto que acontece.

\smallskip\noindent
\textbf{Por que o mercado pode nem fechar (inexistência).} Primeiro, um alívio: o cardápio
separador que achamos ($c^L_\ast = 4$ para o bom motorista) \emph{é} a resposta —
\emph{desde que exista equilíbrio}. O problema mora num caso específico. Chame de $\lambda$ a
fração de motoristas ruins. Se $\lambda$ é \emph{pequena} (quase todo mundo é bom), dá para
oferecer um \emph{pooling} — contrato único para todos, cobertura quase cheia a um preço
médio que, com poucos ruins, sai baixo — e \emph{os dois tipos} ficam melhores do que no
cardápio separador. Isso já \textbf{derruba o separador}: ninguém topa ficar com $25\%$ de
cobertura havendo um pooling superior para ambos. Só que o pooling \textbf{também não se
sustenta}: basta uma seguradora rival lançar um contrato de cobertura parcial e barato que,
por single-crossing, \emph{só} interessa aos bons motoristas; eles migram, e o pooling fica
só com os ruins — e quebra (o preço médio não cobre o custo dos ruins). Moral desconfortável:
com poucos ruins, o separador morre para o pooling e o pooling morre para o
\emph{cream-skimming} — \textbf{nada sobrevive}, e o modelo \emph{pode simplesmente não ter
equilíbrio competitivo}.

\smallskip\noindent
\textbf{E a regulação?} Cuidado para não embaralhar duas falhas. A que mais diretamente
justifica \textbf{ANS e ACA} é a \emph{espiral de seleção adversa} — o \emph{unraveling} do
Exercício~1, em que os bons fogem, o pool encarece e pode colapsar. Mandato de adesão $+$
preço comunitário forçam por lei o pooling que o mercado livre não estabiliza sozinho. A
ligação entre os dois: R-S explica \emph{por que} um pooling não emerge espontaneamente; a
espiral explica \emph{por que} o pooling imposto por lei pode ser socialmente desejável.
\end{empalavras}

\end{enumerate}

\subsection*{Comentário pedagógico (5 passos)}
\begin{gabarito}
\passo{1. Ponto-chave.} No separador R-S, quem se distorce é o \textbf{bom risco} ($L$),
não o ruim. Razão: o $H$ é o tipo \emph{com incentivo a mentir} (quer o seguro barato do
$L$); para impedir isso, encolhe-se a cobertura do $L$ até a IC do $H$ ficar binding. O $H$
sai ileso (cobertura completa, prêmio justo).

\passo{2. Derivação.} Toda a álgebra reduz a $\sqrt{A}+\sqrt{B}=6$ com $A-B=L-c^L_\ast$.
A identidade $(\sqrt{A}-\sqrt{B})(\sqrt{A}+\sqrt{B})=A-B$ evita elevar ao quadrado às cegas.
Na lousa, o atalho honesto é testar $c^L_\ast=4$ (a calibração foi feita para isso); o
caminho fechado fica registrado como quadrática $(c^L)^2-68c^L+256=0$, discriminante
$3600=60^2$.

\passo{3. Armadilha.} (i) Inverter os tipos (``o alto risco deve ficar com menos
seguro''), por intuição moralista — é o oposto. (ii) Tentar achar $c^L_\ast$ impondo
lucro zero ou IR do $L$ — o que define $c^L_\ast$ é a IC do \emph{outro} tipo. (iii)
Aproximar numericamente e perder o fechamento inteiro. (iv) Achar que pooling sempre
sobrevive — se sobrevivesse, R-S não seria interessante.

\passo{4. Extensão.} A inexistência abre a porta para \textbf{regulação}: ANS no Brasil
(Lei 9.656/1998 — rol mínimo, faixas etárias, portabilidade) e ACA nos EUA (2010 —
proibição de \emph{underwriting} + mandato individual + subsídio cruzado) forçam um pooling
que o mercado puro não sustenta, ao custo de o $L$ subsidiar o $H$. Hackmann-Kolstad-Kowalski
(2015, \emph{AER} 105(3): 1030--1066) documentam unraveling parcial em estados que
removeram o mandato — teste empírico de R-S. \textbf{Aula 9} (Spence) traz a via
\emph{descentralizada}: o agente informado sinaliza.

\passo{5. Peer-review (auto-crítica honesta).} Calibração $W=17,\,L=16$ é incomum mas
\emph{deliberada}: ela torna $\sqrt{9}=3$, $\sqrt{16}=4$, $\sqrt{4}=2$ — \emph{todas}
inteiras, ideal para giz. O custo é uma economia onde o sinistro quase zera a riqueza
($W-L=1$); isso não é bug, é \emph{feature} — reforça por que o agente quer tanto seguro.
Os números diferem dos avaliativos (lá $W=100,\,L=56,\,c^L_\ast=8$). A distorção aqui é
$c^L_\ast/L = 4/16 = 25\%$, menos severa que os $14\%$ dos avaliativos e portanto mais
fácil de discutir. Risco residual: o ``teste $c^L_\ast=4$'' parece coelho-da-cartola; por
isso registrei a quadrática fechada como rede de segurança para o aluno cético. Não
derivei a inexistência formalmente — é proposital, está fora do calibre de lousa de 45 min.
\end{gabarito}
\else
\espacoaluno{85mm}
\fi

\hrulefill
\begin{center}\small\textit{%
\ifsolucao Fim dos três exercícios de lousa — Aula 8.%
\else Fim do caderno do aluno — Aula 8. Bom trabalho.\fi}\end{center}

\end{document}
